% GENERATED FILE. Edit canonical lessons, then run npm run generate:books.
% canonical-source-hash: fnv1a64:cdf820b4d958cb5e
% canonical-lessons: FR-C251-resultat, FR-C251-difference, FR-C251-exemple, FR-C251-verite, FR-C251-mensonge

\chapter{Result, Difference, Example, Truth, Lie}
\label{ch:fr-result-difference-example-truth-lie}
\hlchaptermodality{251}

\begin{chapteropening}
\textbf{By the end of this chapter:} \emph{I can say a result, a difference, an example, the truth and a lie.}

Complete the last lesson of chapter 251: I can say a result, a difference, an example, the truth and a lie.
\end{chapteropening}

\section[le résultat]{le résultat — a result}

\label{lesson:FR-C251-resultat}



\begin{quote}
Before the new one: say the French for an elderly person, then the French for a twin.
\end{quote}

\subsection*{You'll want to know: le résultat}
\textbf{le résultat} — \textquotedblleft{}a result\textquotedblright{}.

\begin{cousinweb}
\textbf{résultat} $\leftarrow$ Medieval Latin \emph{resultatum}, from \emph{resultare}, \textquotedblleft{}to spring back\textquotedblright{} — a learned borrowing. (EtymDB 2.1)
\end{cousinweb}

\begin{grammarlens}[title={What you've built}]
One more word for this chapter.
\end{grammarlens}

\subsection*{Your turn}
\begin{itemize}
\raggedright
  \item \emph{Say it:} \emph{le résultat}
  \item \emph{Say it:} \emph{le résultat}, once more
  \item \emph{Say it:} say \emph{le résultat}
  \item \emph{Recall:} say the French for an elderly person, then the French for a twin, then say \emph{le résultat} again
\end{itemize}

\begin{tcolorbox}[breakable,colback=teal!4,colframe=teal!35!black,title={Before you move on}]
What does le résultat mean? (\textquotedblleft{}A result\textquotedblright{}.) Say it once more.
\end{tcolorbox}
\section[la différence]{la différence — a difference}

\label{lesson:FR-C251-difference}



\begin{quote}
Before the new one: say the French for a twin, then the French for a result.
\end{quote}

\subsection*{You'll want to know: la différence}
\textbf{la différence} — \textquotedblleft{}a difference\textquotedblright{}.

\begin{cousinweb}
\textbf{différence} $\leftarrow$ Latin \emph{differentia}, \textquotedblleft{}a difference\textquotedblright{} — the noun beside \emph{différent}. (EtymDB 2.1)
\end{cousinweb}

\begin{grammarlens}[title={What you've built}]
One more word for this chapter.
\end{grammarlens}

\subsection*{Your turn}
\begin{itemize}
\raggedright
  \item \emph{Say it:} \emph{la différence}
  \item \emph{Say it:} \emph{la différence}, once more
  \item \emph{Say it:} say \emph{la différence}
  \item \emph{Recall:} say the French for a twin, then the French for a result, then say \emph{la différence} again
\end{itemize}

\begin{tcolorbox}[breakable,colback=teal!4,colframe=teal!35!black,title={Before you move on}]
What does la différence mean? (\textquotedblleft{}A difference\textquotedblright{}.) Say it once more.
\end{tcolorbox}
\section[l'exemple]{l'exemple — an example}

\label{lesson:FR-C251-exemple}



\begin{quote}
Before the new one: say the French for a result, then the French for a difference.
\end{quote}

\subsection*{You'll want to know: l'exemple}
\textbf{l'exemple} — \textquotedblleft{}an example\textquotedblright{}.

\begin{cousinweb}
\textbf{exemple} $\leftarrow$ Latin \emph{exemplum}, \textquotedblleft{}a sample, a pattern\textquotedblright{} — a learned borrowing. (EtymDB 2.1)
\end{cousinweb}

\begin{grammarlens}[title={What you've built}]
One more word for this chapter.
\end{grammarlens}

\subsection*{Your turn}
\begin{itemize}
\raggedright
  \item \emph{Say it:} \emph{l'exemple}
  \item \emph{Say it:} \emph{l'exemple}, once more
  \item \emph{Say it:} say \emph{l'exemple}
  \item \emph{Recall:} say the French for a result, then the French for a difference, then say \emph{l'exemple} again
\end{itemize}

\begin{tcolorbox}[breakable,colback=teal!4,colframe=teal!35!black,title={Before you move on}]
What does l'exemple mean? (\textquotedblleft{}An example\textquotedblright{}.) Say it once more.
\end{tcolorbox}
\section[la vérité]{la vérité — the truth}

\label{lesson:FR-C251-verite}



\begin{quote}
Before the new one: say the French for a difference, then the French for an example.
\end{quote}

\subsection*{You'll want to know: la vérité}
\textbf{la vérité} — \textquotedblleft{}the truth\textquotedblright{}.

\begin{cousinweb}
\textbf{vérité} $\leftarrow$ Old French \emph{verité} $\leftarrow$ Latin \emph{veritas}, \textquotedblleft{}truth\textquotedblright{}, from \emph{verus}, \textquotedblleft{}true\textquotedblright{} — taken from Latin. (EtymDB 2.1)
\end{cousinweb}

\begin{grammarlens}[title={What you've built}]
One more word for this chapter.
\end{grammarlens}

\subsection*{Your turn}
\begin{itemize}
\raggedright
  \item \emph{Say it:} \emph{la vérité}
  \item \emph{Say it:} \emph{la vérité}, once more
  \item \emph{Say it:} say \emph{la vérité}
  \item \emph{Recall:} say the French for a difference, then the French for an example, then say \emph{la vérité} again
\end{itemize}

\begin{tcolorbox}[breakable,colback=teal!4,colframe=teal!35!black,title={Before you move on}]
What does la vérité mean? (\textquotedblleft{}The truth\textquotedblright{}.) Say it once more.
\end{tcolorbox}
\section[le mensonge]{le mensonge — a lie}

\label{lesson:FR-C251-mensonge}



\begin{quote}
Before the new one: say the French for an example, then the French for the truth.
\end{quote}

\subsection*{You'll want to know: le mensonge}
\textbf{le mensonge} — \textquotedblleft{}a lie\textquotedblright{}.

\begin{cousinweb}
\textbf{mensonge} $\leftarrow$ Old French \emph{mençonge} (EtymDB 2.1), usually traced to Vulgar Latin \emph{*mentionica}, built on Latin \emph{mentiri}, \textquotedblleft{}to lie\textquotedblright{} (Lewis \& Short) — inherited.
\end{cousinweb}

\begin{grammarlens}[title={What you've built}]
One more word for this chapter.
\end{grammarlens}

\subsection*{Your turn}
\begin{itemize}
\raggedright
  \item \emph{Say it:} \emph{le mensonge}
  \item \emph{Say it:} \emph{le mensonge}, once more
  \item \emph{Say it:} say \emph{le mensonge}
  \item \emph{Recall:} say the French for an example, then the French for the truth, then say \emph{le mensonge} again
\end{itemize}

\begin{tcolorbox}[breakable,colback=teal!4,colframe=teal!35!black,title={Before you move on}]
What does le mensonge mean? (\textquotedblleft{}A lie\textquotedblright{}.) Say it once more.
\end{tcolorbox}
